% NR56081C
% Krampfender Patient
% 14.0
% 2013-11-30

:ref:`m-zerebraler-krampfanfall-bestendend`

\Ca{Vitale Bedrohung!}

Taktik Verletzungen vermeiden,
Vitalfunktionen sichern, Abklären ob Status epilepticus
besteht, Klärung des Herganges (Fremdanamnese!),
Differentialdiagnosen abklären}

\begin{itemize}
  -   |TxMassVitMKBes|
  \begin{itemize}
    -   Lagerung: sichern. In der Nachschlafphase stabile
    Seitenlage
    \begin{itemize}
      -   Patient vor Verletzungen schützen, Möbel etc. entfernen oder polstern
      -   NICHT festhalten -- aber sichern, Eigenschutz
      beachten!\footnote{\textbf{Sichern}: Patient darf
      nicht von der Trage oder aus dem Bett fallen, der Kopf
      darf nicht am Boden oder Tischbein anschlagen, etc.}
    \end{itemize}
  \end{itemize}
  -   Besondere Diagnostik:
  \begin{itemize}
    -   Neurocheck inkl. Messung des **Blutzuckers**!
    -   Traumacheck
    -   Nach *jedem Krampanfall*: Einschätzungsblock
    (\FullPrettyRef{m:einschaetzungsblock}) inkl.
    Trauma-Untersuchungen
  \end{itemize}
  -   Reizabgeschirmter Transport
  -   Transportentscheidung: Abt.f. Innere Medizin. (evtl. auch
  Intensivüberwachung)
\end{itemize}

\Customized{
\begin{description}
  \item[\CustAsboe] Beim bestehenden zerebralen Krampfanfall ist die Gabe
  eines Benzodiazepins gem. Algorithmus durch NKA vorgesehen.
 :cite:`Asb921Memo-AL-2011-000001`\ .
\end{description}
}